\section{Packing-preserving sparsification and the final tree}
\label{sec:iteration}


\begin{lemma}[One sparsification step]
\label{lem:ta-sparsification}
For all sufficiently large integers $r$, from an explicit graph $H$
containing $r$ disjoint spanning trees one can construct a subgraph
$H'=(V(H),A)$ and an integer $r'$ such that
\begin{align}
 &r/4\le r'\le r/2,\qquad H'\text{ contains }r'\text{ disjoint spanning trees},
          \label{eq:ta-packing-halving}\\
 &|\delta_{H'}(S)|\le a_r|\delta_H(S)|\quad\text{for every cut},
       \qquad a_r=\tfrac12+3C_s r^{-1/8},
          \label{eq:ta-cut-halving}\\
 &\frac{a_r}{r'/r}\le1+C_1r^{-1/16}
          \label{eq:ta-ratio-loss}
\end{align}
for an absolute $C_1$. The construction is deterministic and takes
polynomial bit time, with the certificate in \Cref{lem:ta-signing}.
\end{lemma}

\begin{proof}
Use \Cref{prop:ta-extraction} and discard excess trees to obtain a union
$B$ of $q$ trees with $(1-Ar^{-1/16})r\le q\le r$.
Increase the cutoff so $q\ge r/2$, and apply
\Cref{lem:ta-frame,lem:ta-signing}. Put $\epsilon_r=C_s r^{-1/8}$.
The two principal blocks of \Cref{eq:ta-half-frame} give
\begin{equation}
 L_A\preceq\tfrac12L_B+\epsilon_rX
      \preceq\tfrac12L_H+\epsilon_rX,
       \qquad \sum_{e\in A}z_ez_e^T\succeq(\tfrac12-\epsilon_r)I.
 \label{eq:ta-two-blocks}
\end{equation}
A centered cut indicator $x$ has $xx^T\in\mathcal N$.
By \Cref{eq:ta-K}, $x^TXx\le3x^TL_Hx$; evaluating the first inequality
on $x$ proves \Cref{eq:ta-cut-halving}.

For packing, fix a partition into $d$ parts. Let $W$ be the span of all
$z_e$ with endpoints in a common part, and let $\Pi$ be the orthogonal
projection onto $W^\perp$. Rank preservation from \Cref{lem:ta-frame}
gives $\dim W\le\sum_{P\in\mathcal P}(|P|-1)=n-d$, hence
$\rank\Pi\ge d-1$. Taking the trace after projection in the second
inequality of \Cref{eq:ta-two-blocks}, and then using
\Cref{eq:ta-frame}, gives
\[
 (\tfrac12-\epsilon_r)(d-1)
       \le\sum_{e\in E_A(\mathcal P)}\|\Pi z_e\|^2
       \le |E_A(\mathcal P)|\frac{1+4/r}{q}.
\]
The integer packing criterion thus supplies at least
\begin{equation}
                     \left\lfloor
                       \frac{q(1/2-\epsilon_r)}{1+4/r}
                     \right\rfloor
 \label{eq:ta-projected-packing}
\end{equation}
spanning trees in $H'$.
Compute its maximum packing number, capped at $\lfloor r/2\rfloor$,
by binary search with the constructive graphic-matroid union algorithm,
and denote the result by $r'$. The lower bound in
\Cref{eq:ta-projected-packing} is already at most $\lfloor r/2\rfloor$,
so capping does not reduce that guarantee. Consequently
\[
 \frac{r'}r\ge\frac12(1-C_2r^{-1/16})
\]
for an absolute $C_2$, absorbing the floor loss $1/r$, the extraction
loss, and $\epsilon_r$. Taking a sufficiently large absolute cutoff
gives \Cref{eq:ta-packing-halving}. Dividing
$a_r=\tfrac12(1+6\epsilon_r)$ by this lower bound and enlarging the
absolute constant proves \Cref{eq:ta-ratio-loss}.
\end{proof}

\begin{proof}[Proof of \Cref{thm:tt-algorithm} for explicit graphs]
First suppose $n\ge2$. Put
$r_{\mathrm{init}}=\max\{1,\lfloor k/2\rfloor\}$.
For a partition into $d\ge2$ parts, summing the $k$-edge-connectivity
bounds on its parts gives $|E_G(\mathcal P)|\ge kd/2$.
Thus \Cref{eq:ta-integer-packing} supplies
$\lfloor k/2\rfloor$ trees, and connectedness supplies one when $k=1$.
Construct this initial packing. Notice that
$r_{\mathrm{init}}\ge k/3$.

Starting with $H_0=G$ and $r_0=r_{\mathrm{init}}$, apply
\Cref{lem:ta-sparsification} whenever the current parameter is at least
the fixed absolute cutoff $R_0$. Write $H_i,r_i$ for the resulting
sequence.
At termination $1\le r_J<R_0$, and every $H_i$ remains connected.
Return any spanning tree of $H_J$.

The ratio estimate gives, for every cut,
\begin{align*}
 |\delta_{H_J}(S)|
 &\le |\delta_G(S)|\prod_{i<J}a_{r_i}\\
 &\le |\delta_G(S)|\,\frac{r_J}{r_{\mathrm{init}}}
                           \prod_{i<J}(1+C_1r_i^{-1/16}).
\end{align*}
Since $r_{i+1}\le r_i/2$, reading the parameters backwards bounds
\[
 \sum_{i<J}r_i^{-1/16}
       \le R_0^{-1/16}\sum_{j\ge0}2^{-j/16}.
\]
Thus the product is bounded by one absolute constant, independently of
the initial packing number. With $r_J<R_0$ and
$r_{\mathrm{init}}\ge k/3$, the displayed estimate implies
\Cref{eq:ta-thin} for a universal $C$. The same calculation covers
$J=0$, with empty product one: in that case
$r_{\mathrm{init}}<R_0$ and any spanning tree suffices.

There are at most $J_{\max}=1+\lceil\log_2(m+1)\rceil$ iterations.
Every iteration consists of deterministic polynomial-time routines,
including the rational signing algorithm of \Cref{thm:rational-signing}.
The arithmetic and precision bounds proved above ensure that all
intermediate encodings and the total work remain polynomial in the
explicit input length.
When $n=1$, return the empty tree before any packing or matrix call.
\end{proof}

\paragraph{Binary multiplicities.}
It remains to reduce a succinct multigraph to an explicit one without
work polynomial in the numerical value of $k$. For $n\ge2$, let
$c_{uv}$ be the binary multiplicity of the unordered pair $uv$ and set
\begin{equation}
 c'_{uv}=\min\{c_{uv},k\},\qquad
 a_{uv}=\left\lfloor\frac{2n^2c'_{uv}}k\right\rfloor.
 \label{eq:ta-binary-reduction}
\end{equation}
Every nontrivial cut still has $c'$-capacity at least $k$: either one
crossing pair already contributes $k$, or no crossing multiplicity was
changed. Rounding loses less than one per crossing pair, of which
there are at most $n^2/4$. Consequently the explicit graph $\widetilde G$
with multiplicities $a_{uv}$ satisfies
\begin{equation}
 |\delta_{\widetilde G}(S)|\ge2n^2-n^2/4\ge n^2,
 \qquad
 |\delta_{\widetilde G}(S)|\le\frac{2n^2}k|\delta_G(S)|.
 \label{eq:ta-binary-cuts}
\end{equation}
Each $a_{uv}\le2n^2$ and there are fewer than $n^2/2$ pairs, so
$|E(\widetilde G)|<n^4$.
Apply the explicit algorithm with connectivity bound $n^2$.
Its output tree is $C/n^2$-thin in $\widetilde G$, and hence
$2C/k$-thin in the original graph by \Cref{eq:ta-binary-cuts}.
Although temporary scaling can make $a_{uv}>c_{uv}$, a tree uses at
most one edge of any parallel pair. Every selected pair has
$c_{uv}>0$, so its tree edge maps to a genuine original edge.
All arithmetic in \Cref{eq:ta-binary-reduction} is polynomial in
$\log k$, the multiplicity encoding lengths, and $n$; the main routine
uses the new parameter $n^2$. Enlarging the universal constant by two
completes \Cref{thm:tt-algorithm} in the binary model.
